\bookchapter{Capítulo 30: Proyección de pantalla, presentaciones y conectividad con el automóvil}

\chapterintro{Por muy grande que sea la pantalla de tu teléfono, sigue siendo de 7 pulgadas; pero tu TV tiene 55 pulgadas y tu proyector 100. “Lanzar” la imagen del teléfono a una pantalla grande es una experiencia de otra dimensión. Este capítulo explica la proyección de pantalla, las presentaciones y la conectividad con el coche; tres escenarios frecuentes que te ayudan a extender las capacidades del teléfono a pantallas más grandes.}

\booksection{I. Los tres campos de la proyección inalámbrica - no los confundas}

La proyección inalámbrica parece “la imagen del teléfono en la TV”, pero hay tres tecnologías subyacentes, no compatibles entre sí:

\begin{center}
\begin{tabularx}{\linewidth}{|X|X|X|}
\hline
\textbf{Tecnología} & \textbf{Quién la usa} & \textbf{Características} \tabularnewline
\hline
AirPlay & iPhone/iPad/Mac & Ecosistema Apple, la experiencia más fluida. Requiere Apple TV o TV compatible con AirPlay \tabularnewline
Miracast & Teléfonos Android & Conexión directa Wi-Fi Direct, sin necesidad de router. La "proyección inalámbrica" de Android usa esto por defecto \tabularnewline
DLNA & Todas las plataformas & Solo envía archivos multimedia (vídeo/música/fotos), no duplica la pantalla. Adecuado para enviar vídeo a la TV \tabularnewline
\hline
\end{tabularx}
\end{center}

\booksubsection{Métodos de operación por plataforma:}

\begin{itemize}
  \item \textbf{Sinceramente, es normal que el móvil empiece a irse lento tras un ano de uso. La causa no necesariamente es que este roto; más a menudo es "acumulacion de basura": archivos basura, procesos en segúndo plano, espacio de almacenamiento llenandose poco a poco.}: Centro de control > Duplicar pantalla > Selecciona tu TV/Apple TV. Si la TV no soporta AirPlay, puedes comprar un Apple TV o un dongle compatible con AirPlay.
  \item \textbf{Android}: Desliza la barra de notificaciones hacia abajo > Busca el botón "proyección inalámbrica" / "proyección inteligente" / "multiventana interactiva" > Selecciona la TV. Requiere que el teléfono y la TV estén en el mismo Wi-Fi.
  \item \textbf{Solución universal}: Si la TV es muy antigua y no soporta ninguna proyección inalámbrica, compra un Chromecast o una Mi Box/Huawei Box y conéctalo al HDMI de la TV para recibir la proyección.
\end{itemize}

\booksection{II. Proyección por cable - más estable, cero latencia}

La proyección inalámbrica es cómoda pero tiene latencia; jugar y hacer presentaciones es mala experiencia. La proyección por cable soluciona esto:

\begin{itemize}
  \item \textbf{Sinceramente, es normal que el móvil empiece a irse lento tras un ano de uso. La causa no necesariamente es que este roto; más a menudo es "acumulacion de basura": archivos basura, procesos en segúndo plano, espacio de almacenamiento llenandose poco a poco.}: Necesitas un cable adaptador de Lightning/USB-C a HDMI (unos 100-200 yuanes). Al conectar, la imagen del teléfono se envía directamente al HDMI de la TV/proyector.
  \item \textbf{Android}: Algunos teléfonos soportan conexión directa USB-C a HDMI (requiere que el teléfono soporte DP Alt Mode, la mayoría de los insignia lo hacen). Compra un cable USB-C a HDMI.
  \item \textbf{Escenario de sala de reuniones}: Si presentas a menudo con el teléfono, te recomiendo llevar siempre un cable a HDMI en la bolsa. Las TV/proyectores de las salas de reuniones casi todos tienen HDMI; conectas y listo, mucho más fiable que pelearse con la proyección inalámbrica.
\end{itemize}

\booksection{III. El teléfono como ordenador - modo escritorio}

Algunos teléfonos Android soportan el “modo escritorio”: al conectar una pantalla, el teléfono genera una interfaz similar al escritorio de un ordenador, operable con ratón y teclado:

\begin{itemize}
  \item \textbf{Samsung DeX}: La función estrella de los Samsung insignia. Al conectar una pantalla, aparece una interfaz de escritorio con ventanas múltiples; la experiencia más cercana a un ordenador real. Con teclado y ratón Bluetooth, puedes viajar sin portátil.
  \item \textbf{Modo escritorio de Huawei}: Los Huawei Mate series/Pura series también tienen una función similar al conectar una pantalla. Soporta ventanas múltiples, gestión de archivos y apps de oficina.
  \item \textbf{Motorola Ready For}: Los Moto insignia también lo soportan; funciones similares a DeX.
  \item \textbf{Casos de uso}: Oficina ligera (responder correos, editar documentos), presentar PPT, ver vídeos en pantalla grande. No adecuado para Excel intensivo o diseño profesional.
\end{itemize}

\booksection{IV. Conectividad con el coche - la forma correcta de usar el teléfono mientras conduces}

Operar el teléfono mientras conduces es ilegal, pero navegar, escuchar música y responder llamadas son necesidades reales. Los sistemas de conectividad del coche te permiten usar las funciones principales del teléfono en la pantalla del coche; el teléfono puede quedarse apartado:

\begin{itemize}
  \item \textbf{Apple CarPlay}: La primera opción para usuarios de iPhone. Conecta por USB o inalámbrico al coche; en la pantalla del coche aparece una interfaz simplificada del iPhone: navegación (Gaode/Baidu Maps), música, teléfono, mensajes, podcasts. Siri permite control por voz sin manos.
  \item \textbf{Android Auto}: La solución equivalente para Android. Funciones similares a CarPlay; soporta Google Maps, Spotify, etc. En China, algunas funciones no están disponibles por restricciones de servicios de Google.
  \item \textbf{Huawei HiCar}: Solución propia de Huawei; alta tasa de soporte en modelos nacionales. Además de navegación y música, ofrece videollamadas Huawei, imágenes de aparcamiento fusionadas y otras funciones especiales.
  \item \textbf{Baidu CarLife}: Solución multiplataforma; funciona tanto en iOS como en Android. Menos funciones pero gran compatibilidad; muchos coches antiguos solo soportan CarLife.
\end{itemize}

\begin{tcbwarn}{Alerta de seguridad}
Sea cual sea la solución de conectividad del coche, mientras conduces deberías usar el control por voz en la medida de lo posible. Configura el destino de navegación con antelación y selecciona la playlist antes de empezar; no operes la pantalla del coche mientras conduces.

\end{tcbwarn}

\booksection{V. Proyección en sala de reuniones - sin estrés}

Las reuniones son el escenario más frecuente de proyección y también donde más problemas surgen. Aquí tienes algunas experiencias prácticas:

\begin{itemize}
  \item \textbf{Llega temprano para probar}: Nunca esperes a que empiece la reunión para conectar la proyección. Llega 10 minutos antes y prueba el método de conexión. La probabilidad de fallo de la proyección Wi-Fi es mucho mayor de lo que crees.
  \item \textbf{Prepara una solución por cable como Plan B}: Lleva un cable a HDMI en la bolsa. Si la proyección inalámbrica no funciona, el cable salva la vida.
  \item \textbf{Desactiva la vista previa de notificaciones}: Al proyectar, tus mensajes de WeChat aparecerán en la pantalla grande. Antes de proyectar, desactiva "vista previa del contenido de notificaciones en pantalla de bloqueo" en los ajustes, o activa el "modo presentación" (algunos teléfonos lo soportan; oculta automáticamente las notificaciones durante la proyección).
  \item \textbf{Usa "presentación de documentos" en lugar de "operar el teléfono"}: Si vas a mostrar un PPT, guárdalo en el teléfono y ábrelo a pantalla completa con WPS en lugar de proyectar toda la pantalla del teléfono. Es más profesional y fluido.
\end{itemize}

\begin{tcbpractice}{Práctica: completa una proyección de teléfono a TV}
\begin{enumerate}
  \item Cambio de plataforma cruzado (Android a iPhone o viceversa)
  \item Abre una foto o un vídeo en el teléfono.
  \item Intenta la proyección inalámbrica (iPhone usa duplicación de pantalla, Android usa el botón de proyección inalámbrica).
  \item Si falla, comprueba si la TV soporta el protocolo correspondiente o prueba DLNA (botón "proyectar" en la app de galería/vídeo).
  \item Tras proyectar con éxito, prueba a hacer otras operaciones en el teléfono para ver si la pantalla grande se sincroniza.
\end{enumerate}

\end{tcbpractice}

\begin{tcbsummary}{Resumen del capítulo}
\begin{itemize}
  \item Tres campos de proyección inalámbrica: AirPlay (Apple), Miracast (Android), DLNA (multimedia push multiplataforma)
  \item La proyección por cable es más estable con cero latencia; indispensable para sala de reuniones con cable a HDMI
  \item Samsung DeX / modo escritorio de Huawei convierten el teléfono en un ordenador de oficina ligera
  \item Conectividad del coche: cuatro soluciones CarPlay / Android Auto / HiCar / CarLife
  \item Reglas de oro de proyección en reuniones: prueba con antelación, prepara un Plan B, desactiva la vista previa de notificaciones
\end{itemize}

\end{tcbsummary}
